\documentclass{article}

\usepackage{amsmath,amssymb}
\usepackage{graphicx}

\usepackage{pgfpages}

\title{Plasma Turbulence}
\author{Jayneel Parikh}
\date{September 1, 2026}

\begin{document}


%========================
% Slide 1: Title
%========================
\section{}
\begin{itemize}
\item This talk is about plasma turbulence and why it is the dominant limitation to magnetic confinement fusion.
\item The focus is not on turbulence modeling itself, but on how we interrogate experimental fluctuation data.
\item In particular, I will discuss nonlinear, directional correlation methods developed to go beyond standard diagnostics.
\end{itemize}

%========================
% Slide 2: Historical Timeline
%========================
\section{}
\begin{itemize}
\item Fusion research originates from early nuclear physics, beginning with Rutherford’s first fusion reaction in 1919.
\item By the 1920s, mass defect and nuclear binding energy explained how fusion could release enormous energy.
\item Eddington proposed that stellar energy originates from hydrogen fusion.
\item In 1938, Bethe and Weizsäcker established the proton--proton chain and CNO cycle.
\item In the 1950s, fusion research accelerated through weapons programs and the invention of toroidal confinement concepts.
\item The 1968 Soviet T--3 tokamak demonstrated dramatically improved confinement, launching modern tokamak research.
\item Since the 1990s, the primary challenge has not been ignition physics, but turbulence and confinement degradation.
\item Modern devices like ITER and EAST are fundamentally limited by turbulence-driven transport.
\end{itemize}
\section{}
%========================
% Slide 3: Plasma and Tokamaks
%========================
\begin{itemize}
\item Plasma is the fourth state of matter, consisting of charged particles interacting through Coulomb forces.
\item Unlike neutral fluids, plasma dynamics are strongly constrained by magnetic fields.
\item A tokamak confines plasma magnetically in a toroidal geometry to avoid material walls.
\item Charged particles spiral along magnetic field lines while slowly drifting across them.
\item Toroidal and poloidal magnetic fields combine to form closed helical particle orbits.
\item Ideal confinement would restrict transport across magnetic flux surfaces.
\end{itemize}

\section{}
%========================
% Slide 4: The Confinement Problem — Turbulence
%========================
\begin{itemize}
\item Experimental transport levels exceed classical and neoclassical predictions by one to two orders of magnitude.
\item This discrepancy is known as anomalous transport.
\item Turbulence arises from drift-type instabilities driven by pressure gradients and electric fields.
\item Plasma turbulence is anisotropic, intermittent, and non-Gaussian.
\item Zonal flows represent a form of self-organization analogous to atmospheric bands on Jupiter.
\item Modern transport analysis reveals both slow diffusive transport and fast nonlocal transport events.
\item These fast events resemble Lévy flights rather than standard random walks.
\end{itemize}

\section{}

%========================
% Slide 5: Experimental Setup
%========================
\begin{itemize}
\item Experimental turbulence measurements are typically obtained using Langmuir probes.
\item Floating potential fluctuations provide access to local electrostatic turbulence.
\item Multiple spatially separated probes allow the study of radial and poloidal correlations.
\item The experiments referenced here follow setups used by Bsharat et al.\ and Alipour et al.
\item These measurements provide time series suitable for nonlinear statistical analysis.
\end{itemize}

\section{}

%========================
% Slide 6: Interrogating the Data
%========================
\begin{itemize}
\item Conventional analysis tools include linear cross-correlation, FFTs, and coherence.
\item These methods assume Gaussian statistics and waveform coherence.
\item Plasma turbulence violates these assumptions due to shear, intermittency, and deformation.
\item Linear correlation is symmetric and cannot identify causal or directional relationships.
\item This motivates the use of nonlinear, phase-space-based techniques.
\end{itemize}

\section{}

%========================
% Slide 7: Ding vs Transfer Entropy
%========================
\begin{itemize}
\item Ding et al.\ introduced a nonlinear cross-correlation method based on conditional dispersion.
\item The method reconstructs local dynamics using delay-coordinate embedding.
\item An asymmetry coefficient $S_{xy}$ quantifies directional coupling between signals.
\item If $S_{xy} > 0$, energy or information flows from $x$ to $y$.
\item van Milligen et al.\ use Transfer Entropy, an information-theoretic measure of causality.
\item Transfer Entropy quantifies how much the past of one signal predicts the future of another.
\item Both methods aim to extract directional information flow in nonlinear systems.
\end{itemize}

\section{}

%========================
% Slide 8: Core Aspects of NLCC
%========================
\begin{itemize}
\item Time series are embedded into phase space using delay vectors.
\item The embedding dimension $M$ and delay $m$ are chosen using autocorrelation times.
\item Conditional dispersion measures how spread out one attractor is when the other is locally constrained.
\item A neighborhood size $\epsilon$ is swept to construct dispersion curves.
\item The nonlinear correlation strength $g_{xy}$ is extracted from $\epsilon_{90}$.
\item Directionality is determined using the asymmetry coefficient $S_{xy}$.
\end{itemize}

\section{}

%========================
% Slide 9: Applying the Methods
%========================
\begin{itemize}
\item The algorithm is first validated using coupled Van der Pol oscillators.
\item The direction of coupling is known a priori in the synthetic system.
\item This provides a controlled benchmark for the nonlinear method.
\item The method is then applied to experimental HT--7 tokamak data.
\item Results are compared directly against Ding’s original MATLAB implementation.
\end{itemize}
\section{}

%========================
% Slide 10: Results
%========================
\begin{itemize}
\item Dispersion curves for the Van der Pol system recover the correct coupling direction.
\item HT--7 dispersion curves match previously published results.
\item Nonlinear correlation lengths are significantly larger than linear correlation lengths.
\item Direction reversals are observed across confinement transitions.
\item These results support the physical interpretation of inward and outward turbulence propagation.
\end{itemize}
\section{}

%========================
% Slide 11: Future Research
%========================
\begin{itemize}
\item Fully reproduce Ding’s algorithm in a modern, open-source language.
\item Validate results against synthetic and experimental data sets.
\item Establish mathematical links between NLCC, Transfer Entropy, and CTRW models.
\item Interpret experimental results within a stochastic transport framework.
\item Quantify the relative contribution of fast and slow transport channels.
\item Ultimately connect directional turbulence propagation to confinement degradation.
\end{itemize}

\end{document}
